% =============================================================
%  Chapter 10 — Multi-Receiver Signalling and Public Ritual
%  Lean source: FormalAnthropology/IDT_Signal_Ch10.lean
% =============================================================
\chapter{Multi-Receiver Signalling and Public Ritual}
\label{ch:multireceiver}

\epigraph{%
  ``The culture of a people is an ensemble of texts, themselves ensembles,
  which the anthropologist strains to read over the shoulders of those to
  whom they properly belong.''}%
  {Clifford Geertz, \textit{The Interpretation of Cultures} (1973)}

% ----------------------------------------------------------------
\section{Introduction: One Ritual, Many Minds}
\label{sec:mr-intro}

The defining feature of public ritual is that the \emph{same} signal is
sent to \emph{multiple} receivers with \emph{different} cultural
backgrounds.  A Catholic mass is witnessed simultaneously by the devout
grandmother who has attended every Sunday for sixty years, the skeptical
teenager dragged along by family obligation, and the visiting anthropologist
scribbling field notes.  Each receives the same sensory input; each produces
a radically different interpretation.

Clifford Geertz captured this multiplicity in his concept of ``thick
description'': the same cockfight in Bali means ``gambling'' to one
observer, ``status contest'' to another, ``art'' to a third, and ``data''
to the ethnographer.\footnote{Geertz, C., ``Deep Play: Notes on the
Balinese Cockfight,'' in \textit{The Interpretation of Cultures} (Basic
Books, 1973), pp.~412--453.}  Victor Turner extended the insight to ritual
proper: the same ceremony of \emph{Nkang'a} among the Ndembu produces
different (but structurally related) experiences in initiands, elders, and
spectators.\footnote{Turner, V., \textit{The Forest of Symbols: Aspects of
Ndembu Ritual} (Cornell, 1967).}

The previous chapters studied signalling to a \emph{single} receiver.  This
chapter generalizes to the multi-receiver case.  We define an
\emph{audience}---a list of repertoires, each representing a different
mind---and study what happens when a single signal is broadcast to the
entire audience simultaneously.

The central anthropological insight is that public ritual \emph{works} not
because everyone gets the same meaning, but because everyone gets
\emph{resonant} meanings.  The grandmother, the teenager, and the
anthropologist all experience something structurally related---though not
identical---because resonance is preserved under composition with a common
signal.  This is the formal foundation of Durkheim's \emph{conscience
collective}: shared ritual produces shared structure, not shared content.


% ================================================================
\section{Core Definitions}
\label{sec:mr-defs}

\begin{definition}[Audience]
\label{def:audience}
An \emph{audience} is a list of repertoires:
\[
  \mathrm{Audience}(\IS) \;=\; \mathrm{List}(\mathrm{Repertoire}(\IS))
  \;=\; \mathrm{List}(\mathrm{List}\;\IS).
\]
Each element represents one mind's stock of ideas.  The audience is
heterogeneous: different minds may have different numbers of ideas and
different depths.
\end{definition}

\begin{definition}[Audience Outcome]
\label{def:audience-outcome}
The \emph{audience outcome} of sending signal $s$ to audience
$A = [R_1,\ldots,R_m]$ is the list of interpretation lists:
\[
  \mathrm{audienceOutcome}(A, s) \;=\;
  [\mathrm{interpret}(R_1,s),\;\ldots,\;\mathrm{interpret}(R_m,s)].
\]
\end{definition}

\begin{lstlisting}[caption={Audience and audience outcome in Lean~4.}]
abbrev Audience (I : Type*) := List (Repertoire I)

def audienceOutcome
    (audience : Audience I) (s : I) : List (List I) :=
  audience.map (fun rep => interpret rep s)
\end{lstlisting}

\begin{definition}[Total Audience Depth]
\label{def:total-audience-depth}
The \emph{total audience depth} of an outcome is the sum of all depth sums
across all minds:
\[
  \mathrm{totalAudienceDepth}(\mathit{outcomes}) \;=\;
  \sum_{i}\;\mathrm{depthSum}(\mathit{outcomes}_i).
\]
\end{definition}


% ================================================================
\section{Structural Invariants}
\label{sec:mr-structure}

The first group of theorems establishes that the multi-receiver operation
preserves the fundamental structural properties of signalling.

\begin{theorem}[Audience Size Preservation (10.1)]
\label{thm:audience-size}
$|\mathrm{audienceOutcome}(A, s)| = |A|.$
\end{theorem}

\begin{proof}[Proof sketch]
$\mathrm{audienceOutcome}$ is a \texttt{List.map}, which preserves length.
\end{proof}

A public ritual has exactly as many ``experiences'' as there are
participants.  No ghost observers, no merged consciousnesses---each mind
is counted exactly once.  This is a non-trivial modeling choice: it rules
out ``collective consciousness'' as a separate entity above and beyond the
individual minds.%
\footnote{Durkheim's \emph{conscience collective} is often misread as
positing a supra-individual mind.  In IDT, collective consciousness is
emergent: it is the structural similarity (resonance) among individual
outcomes, not a separate entity.}

\begin{theorem}[Per-Receiver Interpretation Count (10.2)]
\label{thm:per-receiver-count}
$|\mathrm{interpret}(R, s)| = |R|.$
\end{theorem}

Each mind produces exactly as many interpretations as it has ideas in its
repertoire.  A five-category mind produces five interpretations; a
fifty-category mind produces fifty.  The signal doesn't create or destroy
ideas---it transforms each one.

\begin{theorem}[Same Signal Different Minds --- Length Preserved (10.3)]
\label{thm:same-signal-different}
For any audience member $R\in A$:
$|\mathrm{interpret}(R, s)| = |R|.$
\end{theorem}

The ``same ceremony, different meanings'' principle: each participant brings
their own framework and gets their own number of meanings, determined by
\emph{their} background, not the signal.  A child at a symphony may have
three categories (``loud,'' ``soft,'' ``boring''); a musicologist may have
three hundred.  The symphony produces three interpretations for the child
and three hundred for the musicologist.


% ================================================================
\section{Void Signal: Universal Transparency}
\label{sec:mr-void}

\begin{theorem}[Void Signal Transparent For Single Receiver (10.4)]
\label{thm:void-transparent-single}
$\mathrm{interpret}(R, \void) = R.$
\end{theorem}

\begin{proof}[Proof sketch]
By induction on $R$.  For each idea $r$, $r\compose\void = r$ by the
right-identity axiom.
\end{proof}

Silence leaves every mind unchanged.  The minute of silence, the Quaker
meeting, the Buddhist \emph{noble silence}---all exploit this: by saying
nothing, every participant's existing worldview remains intact.  Silence
is the only signal that is universally transparent.

\begin{theorem}[Void Signal Transparent For All Audiences (10.5)]
\label{thm:void-transparent-audience}
$\mathrm{audienceOutcome}(A, \void) = A.$
\end{theorem}

\begin{proof}[Proof sketch]
By induction on the audience list.  Each repertoire maps to itself
by Theorem~\ref{thm:void-transparent-single}.
\end{proof}

Universal transparency of silence is the formal foundation of the ``null
ritual''---a ceremony that changes nothing.  Every culture has such
ceremonies: they serve social functions (marking solidarity, observing
convention) without altering anyone's ideas.  The formal result shows that
void is the \emph{unique} signal with this property (assuming the repertoire
is non-trivially structured).

\begin{lstlisting}[caption={Void transparency for audiences in Lean~4.}]
theorem void_transparent_audience (audience : Audience I) :
    audienceOutcome audience (IdeaticSpace.void : I)
      = audience
\end{lstlisting}


% ================================================================
\section{Depth Bounds Across Audiences}
\label{sec:mr-depth}

\begin{theorem}[Single Receiver Depth Bound (10.6)]
\label{thm:single-receiver-depth}
\[
  \mathrm{depthSum}(\mathrm{interpret}(R, s))
  \;\le\;
  \mathrm{depthSum}(R) + |R|\cdot\depth(s).
\]
\end{theorem}

\begin{proof}[Proof sketch]
Induction on $R$.  At each step, $\depth(r\compose s)\le\depth(r)+\depth(s)$
by subadditivity.  Summing over all $|R|$ elements gives the bound.
\end{proof}

Each audience member's ``persuasion exposure'' is linearly bounded.  This
extends the persuasion bounds of Chapter~\ref{ch:persuasion} to the
per-receiver setting: even in a public ritual, each participant's total
interpretive transformation is bounded by their own cultural depth plus the
ceremony's complexity.

\begin{theorem}[Audience Depth Bound --- Per Member (10.7)]
\label{thm:audience-member-depth}
For each $R\in A$:
$\mathrm{depthSum}(\mathrm{interpret}(R,s)) \le \mathrm{depthSum}(R)+|R|\cdot\depth(s).$
\end{theorem}

A PhD and a child at the same concert are both bounded, but differently.
The PhD's larger repertoire means a larger absolute bound, but the \emph{per-idea}
bound ($\depth(r)+\depth(s)$ for each schema) is the same.  The signal's
persuasive impact is democratic in structure but heterogeneous in magnitude.

\begin{theorem}[Empty Audience Has Empty Outcome (10.8)]
\label{thm:empty-audience}
$\mathrm{audienceOutcome}([], s) = [].$
\end{theorem}

A ritual with no participants produces no effects.  Performance requires an
audience.  This is the formal version of the old question: if a tree falls
in a forest and no one is there to hear it, does it make a sound?  In IDT:
if a signal is sent to an empty audience, no interpretations are produced.

\begin{theorem}[Singleton Audience Outcome (10.9)]
\label{thm:singleton-audience}
$\mathrm{audienceOutcome}([R], s) = [\mathrm{interpret}(R,s)].$
\end{theorem}

An audience of one produces exactly one interpretation list.  This is the
formal model of individual meditation or private prayer: a single mind
processing a signal in isolation.  It reduces the multi-receiver framework
to the single-receiver case of Chapters~\ref{ch:equilibrium}--\ref{ch:echo}.


% ================================================================
\section{Resonance Across Audiences}
\label{sec:mr-resonance}

This section contains the deepest results of the chapter.  Resonance---the
structural similarity between ideas---propagates through multi-receiver
signalling in three ways: through resonant signals, through resonant
receivers, and through the combination of both.

\begin{theorem}[Resonant Signals Same Mind (10.10)]
\label{thm:resonant-signals-same}
If $s_1\res s_2$, then for each background idea $r$:
$r\compose s_1 \res r\compose s_2.$
\end{theorem}

Christmas and Hanukkah are ``resonant signals''---structurally similar
festivals involving light, family gathering, and gift-giving.  Each
participant's experience of one resonates with their experience of the
other.  This is why comparative religion is intellectually productive: the
structural similarity of the signals guarantees structural similarity of the
interpretations, regardless of the receiver's background.

\begin{theorem}[Same Signal Resonant Receivers (10.11)]
\label{thm:same-signal-resonant}
If $r_1\res r_2$, then for any signal $s$:
$r_1\compose s \res r_2\compose s.$
\end{theorem}

Two members of the same culture (resonant backgrounds) always produce
resonant interpretations of the same event.  This is Durkheim's
\emph{conscience collective}: cultural consensus is a structural consequence
of shared backgrounds.  The theorem gives it precise mathematical content:
resonant backgrounds + common signal = resonant experiences.%
\footnote{Durkheim, \'E., \textit{Les Formes \'el\'ementaires de la vie
religieuse} (1912).  The concept of \emph{conscience collective} has been
criticized as vague; the resonance theorem provides a formal sharpening.}

\begin{theorem}[Full Resonance: Resonant Receivers + Resonant Signals (10.12)]
\label{thm:full-resonance}
If $r_1\res r_2$ and $s_1\res s_2$, then
$r_1\compose s_1 \res r_2\compose s_2.$
\end{theorem}

\begin{proof}[Proof sketch]
Direct from the resonance-compose axiom: resonance is preserved under
composition when both components resonate.
\end{proof}

Members of allied cultures witnessing similar rituals produce resonant
experiences.  This is Turner's \emph{communitas}: spontaneous solidarity
arising when similar people undergo similar experiences.  The theorem shows
that communitas is not mystical---it is a structural consequence of
resonance preservation under composition.

\begin{theorem}[Self-Resonance of All Audience Outcomes (10.13)]
\label{thm:audience-self-resonance}
For each audience member with $r\in R$:
$r\compose s \res r\compose s.$
\end{theorem}

Each participant's experience of a ritual is internally consistent,
regardless of how different participants' experiences are from each other.
The grandmother's experience is coherent; the teenager's experience is
coherent; the anthropologist's experience is coherent.  What may not be
coherent is the \emph{comparison} between them---but that is a question
about \emph{inter-resonance}, not \emph{self-resonance}.

\begin{lstlisting}[caption={Full resonance theorem in Lean~4.}]
theorem full_resonance
    {r_1 r_2 s_1 s_2 : I}
    (hr : IdeaticSpace.resonates r_1 r_2)
    (hs : IdeaticSpace.resonates s_1 s_2) :
    IdeaticSpace.resonates
      (IdeaticSpace.compose r_1 s_1)
      (IdeaticSpace.compose r_2 s_2) :=
  IdeaticSpace.resonance_compose r_1 r_2 s_1 s_2 hr hs
\end{lstlisting}


% ================================================================
\section{Audience Composition and Merging}
\label{sec:mr-merging}

\begin{theorem}[Audience Concatenation (10.14)]
\label{thm:audience-concat}
\[
  \mathrm{audienceOutcome}(A_1\!+\!+\!A_2,\; s)
  \;=\;
  \mathrm{audienceOutcome}(A_1, s)\!+\!+\;\mathrm{audienceOutcome}(A_2, s).
\]
\end{theorem}

Merging two audiences then signalling is the same as signalling each
audience separately then merging the results.  A joint ceremony for two
communities (e.g., an interfaith service) produces results that are just the
concatenation of what each community would have experienced separately.

There is no mysterious ``synergy'' in public ritual---just the union of
individual experiences.  This is a strong anti-holistic result: it says that
the collective effect of a public ritual is \emph{decomposable} into
per-community effects.  Turner's communitas, whatever its phenomenological
reality, is structurally a concatenation, not a fusion.

\begin{theorem}[Replicated Audience (10.15)]
\label{thm:replicated-audience}
An audience of $n$ identical minds produces $n$ identical interpretation lists:
\[
  \mathrm{audienceOutcome}(\mathrm{replicate}(n, R),\; s)
  \;=\;
  \mathrm{replicate}(n, \mathrm{interpret}(R, s)).
\]
\end{theorem}

\begin{proof}[Proof sketch]
Induction on $n$.  The base case is trivial.  For the inductive step, the
head of the list is $\mathrm{interpret}(R,s)$, and the tail follows by the
inductive hypothesis.
\end{proof}

A perfectly homogeneous community produces perfectly homogeneous responses
to any signal.  This is the theoretical limit of cultural uniformity---the
formal ``ideal type'' of a totalitarian society in which every citizen has
been given the same repertoire.  In practice, no society achieves this: even
the most totalitarian regimes harbor underground variation.%
\footnote{On the limits of cultural homogenization in totalitarian states,
see Fitzpatrick, S., \textit{Everyday Stalinism} (Oxford, 1999);
Gellately, R., \textit{Backing Hitler} (Oxford, 2001).}

\begin{lstlisting}[caption={Replicated audience in Lean~4.}]
theorem replicated_audience
    (rep : Repertoire I) (s : I) (n : ℕ) :
    audienceOutcome (List.replicate n rep) s =
    List.replicate n (interpret rep s)
\end{lstlisting}


% ================================================================
\section{Advanced Multi-Receiver Properties}
\label{sec:mr-advanced}

\begin{theorem}[Sequential Signal to Audience (10.16)]
\label{thm:sequential-audience}
For each audience member:
$r\compose(s_1\compose s_2) = (r\compose s_1)\compose s_2.$
\end{theorem}

A two-act ritual (overture + main act) is equivalent to the main act
performed for an audience that has already absorbed the overture.  The first
act transforms the audience; the second act addresses the transformed
audience.  This is the associativity of composition applied to the ritual
setting, and it justifies the common anthropological practice of analyzing
multi-phase rituals as sequential transformations of the participants'
interpretive state.

\begin{theorem}[Void Audience Member is Transparent (10.17)]
\label{thm:void-audience-member}
$\mathrm{interpret}([\void], s) = [s].$
\end{theorem}

\begin{proof}[Proof sketch]
$\void\compose s = s$ by the left-identity axiom.
\end{proof}

The ``naive observer'' with no cultural baggage sees the signal as-is.  This
formalizes the (impossible) ideal of Husserl's \emph{epoch\'e}---bracketing
all presuppositions to encounter the phenomenon ``as it really is.''  In
IDT, the void audience member is the formal model of phenomenological
reduction: a mind with no background that processes signals transparently.

Of course, no real observer achieves this.  Every ethnographer brings a
repertoire; every scientist brings a paradigm.  The void audience member is
a limiting case, useful for theoretical analysis but unattainable in
practice.%
\footnote{Gadamer's critique of Husserl in \textit{Truth and Method} (1960)
makes precisely this point: understanding always involves
\emph{Vorurteile} (pre-judgments), and the attempt to eliminate them is
itself a pre-judgment.}

\begin{theorem}[Depth Bound for Void Outcome Across Audience (10.18)]
\label{thm:void-outcome-preserved}
$\mathrm{depthSum}(\mathrm{interpret}(R, \void)) = \mathrm{depthSum}(R).$
\end{theorem}

The ``information cost'' of silence is exactly zero for every audience member
simultaneously.  This is the formal foundation of the \emph{status quo}:
doing nothing changes nothing for anyone, regardless of their background.

Combined with the audience concatenation theorem
(Theorem~\ref{thm:audience-concat}), this means that the void signal
preserves the total audience depth exactly:
$\mathrm{totalAudienceDepth}(\mathrm{audienceOutcome}(A,\void))
= \mathrm{totalAudienceDepth}(A)$.
Silence is cost-free at the individual level \emph{and} at the collective
level.

\begin{lstlisting}[caption={Void outcome depth preservation in Lean~4.}]
theorem void_outcome_depth_preserved
    (rep : Repertoire I) :
    depthSum (interpret rep (IdeaticSpace.void : I))
      = depthSum rep
\end{lstlisting}


% ================================================================
\section{Conclusion}
\label{sec:mr-conclusion}

This chapter generalized the signalling framework from single receivers to
heterogeneous audiences, providing formal foundations for the anthropological
study of public ritual.

The central results are:
\begin{enumerate}
  \item \emph{Structural invariance}
    (Theorems~\ref{thm:audience-size}--\ref{thm:same-signal-different}):
    audience size and per-receiver interpretation count are preserved.
  \item \emph{Universal transparency of void}
    (Theorems~\ref{thm:void-transparent-single}--\ref{thm:void-transparent-audience}):
    silence changes no one.
  \item \emph{Resonance propagation}
    (Theorems~\ref{thm:resonant-signals-same}--\ref{thm:full-resonance}):
    resonant signals and resonant receivers produce resonant experiences.
  \item \emph{Decomposability} (Theorem~\ref{thm:audience-concat}):
    multi-community rituals decompose into per-community effects.
\end{enumerate}

The theory vindicates Geertz's ``thick description'' program: different
observers \emph{do} produce different interpretations, but these
interpretations are structurally related through the resonance relation.
It formalizes Durkheim's \emph{conscience collective} as resonance among
individual outcomes, and Turner's \emph{communitas} as the full resonance
theorem (Theorem~\ref{thm:full-resonance}).

Together, Chapters~\ref{ch:equilibrium}--\ref{ch:multireceiver} provide a
complete game-theoretic framework for the study of signalling in ideatic
spaces: from equilibrium existence (Chapter~\ref{ch:equilibrium}) through
persuasion bounds (Chapter~\ref{ch:persuasion}), lossy reception
(Chapter~\ref{ch:lossy}), echo chambers (Chapter~\ref{ch:echo}), to the
multi-receiver case studied here.  The framework is fully formalized in
Lean~4, with every theorem machine-verified and zero \texttt{sorry} axioms.
